知识地图给出了各部分在体系中的位置，学习路线要解决的是更具体的问题：以当前目标和已有知识为起点，先走哪一段，什么时候回补，在哪里汇入完整框架。

这套体系按数学依赖组织，却不要求每位读者都从第一页线性读到最后一页。真正需要保持的是概念之间的先后关系：若一个结论依赖尚未理解的对象，就沿知识地图向前回补；若已经能解释公式的对象、条件与失败方式，就不必为了“完整”重复停留。

下面给出四条入口。它们不是四份互不相干的书单，而是对同一张知识网络的四种走法。

\subsection{路线一：建立完整框架}

若目前能代入公式，却常说不清公式为什么成立，建议从数学语言开始，依次经过微积分、线性代数、概率论与数理统计。离散数学、随机过程、信息论、数值分析和优化随后补上结构、时间依赖、信息尺度与计算可靠性。矩阵微积分把这些工具接到复合模型的梯度上，实分析支撑则负责解释收敛、连续与极限交换的条件。

这条路线不要求在第一次阅读时完成所有严格证明。第一次经过一个概念时，先抓住对象、例子和结论；当后续推导真的调用它，再回来补证明中的技术条件。严格性不是在开头设置的一道门，而是在使用过程中逐步偿还的债务。

\subsection{路线二：为机器学习补数学}

若目标是理解模型而不是只会调用接口，可以先读线性代数中的投影、最小二乘、谱与 SVD，再读多元微积分、矩阵微积分、概率论、数理统计和凸优化。随后进入机器学习，沿``学习问题与可信评价---线性模型---概率模型---近似推断---学习理论''建立主线。

遇到深度模型时，再进入深度学习中的函数复合、反向传播、训练动力学、卷积、注意力与变分生成模型。这里有三个不能互相替代的问题：模型能否表示目标函数，训练算法能否找到合适参数，有限样本能否支持对新数据的结论。表示、优化和泛化必须分别寻找证据。

在凸优化与深度学习训练之间，可以先经过``非凸与随机优化：能下降不等于找到全局解''。它把全局最优证书降为驻点、负曲率或 PL 条件下的对应结论，并解释 SGD、momentum 与 Adam 的理论保证为何弱于它们的更新式外观。

若最终问题涉及治疗、政策、产品干预或机制解释，应在可信评价之后进入``因果推断：预测世界不等于改变世界''。预测模型学习 $P(Y\mid X)$，因果分析则要额外说明干预目标、处理分配机制与可识别性；测试集准确率不能代替这层论证。

\subsection{路线三：概率、统计与序贯决策}

若主要关心不确定性与决策，可以从概率空间、条件概率、随机变量和期望开始，再读联合分布、极限定理、随机过程与信息论。数理统计把``已知分布算概率''转成``给定数据判断未知机制''，机器学习中的概率模型和近似推断进一步处理潜变量与不可解析 posterior。

若问题最终必须采取行动，应在点估计、区间与检验之后继续读``统计决策与可信预测：模型给出概率之后怎样行动''。它把 posterior、频率风险、校准、预测集合、拒绝成本和容量约束放在同一条链上，防止把“概率算得出来”误读成“行动已经确定”。

当推导开始依赖零概率条件事件、混合分布、换测度或期望与极限交换时，再回到实分析支撑中的``测度、积分与条件期望：为现代概率回填地基''。这一步不必挡在初等概率之前，但应在系统使用 martingale、重要性采样或连续状态 Bellman 方程之前补上。

强化学习还需要 Markov 链、条件期望、随机逼近和动态规划。最优控制先研究模型已知时怎样解值函数或最优轨线，强化学习再研究只能观察样本转移时怎样逼近同一个 Bellman 结构。

\subsection{路线四：按故障定向查缺}

当问题来自一条具体公式或一次失败的计算，可以从症状反向定位：

\begin{itemize}
\tightlist
\item
  矩阵乘不起来、转置位置混乱：回到线性映射、张量轴和矩阵微积分；
\item
  梯度与数值差分不一致：检查微分、链式法则、反向传播与浮点误差；
\item
  训练损失下降而验证表现变差：检查经验风险、数据划分、容量与分布偏移；
\item
  posterior 写得出却无法积分：进入 Monte Carlo、MCMC、变分推断与序贯 Monte Carlo；
\item
  迭代震荡、发散或对初值极敏感：检查条件数、步长、凸性、稳定性与动力系统；
\item
  强化学习更新式看似合理却不收敛：区分采样误差、bootstrap、离策略分布和函数逼近。
\end{itemize}

定向查缺不等于只记住修复技巧。修复之后应能回答：错误属于问题定义、数学模型、优化算法、有限精度实现，还是统计评价？只有定位到这一层，经验才可能迁移到下一次失败。

\subsection{什么时候可以继续向后读}

判断是否掌握一节，不必以能否背诵全部公式为标准。更可靠的信号是：能说清公式中的对象和维度，能指出结论依赖的条件，能用一个小例子重建关键推导，并能给出至少一种失效情形。若只能复述结论却无法制造反例，理解通常还停留在名称层面；若能解释为什么某个看似合理的替换会破坏前提，就已经具备继续迁移的基础。
